\documentclass[10pt,a4paper]{amsart}
\usepackage[margin=19mm]{geometry}
\usepackage{amsmath,amssymb,mathtools}
\usepackage[scr=rsfs,cal=euler]{mathalfa}
\usepackage{xfrac}
\usepackage[unicode,hidelinks,bookmarks=false]{hyperref}
\newcommand{\ie}{i.e.\ }
\newcommand{\cf}{cf.\ }
\newcommand{\resp}{respectively\ }
\newcommand{\Subsection}{\subsection}
\providecommand{\nobreakdash}{\nobreak\hbox{-}}
\setlength{\emergencystretch}{2em}
\multlinegap0pt
\usepackage{xcolor}
\usepackage{amssymb}
\usepackage{enumerate}
\usepackage[shortlabels]{enumitem}
\usepackage{tikz}
\newtheorem{proposition}{Proposition}
\newtheorem{definition}{Definition}
\newtheorem{lemma}{Lemma}
\newtheorem{theorem}{Theorem}
\title{Asymptotics of Bounded Lecture-Hall Tableaux}
\author{David Keating, Zhongyang Li, Istvan Prause}
\date{Source: arXiv:2309.15235v3 (CC BY 4.0)}
\begin{document}
\maketitle

\noindent\emph{Source and licence.} Asymptotics of Bounded Lecture-Hall Tableaux. Version arXiv:2309.15235v3 (CC BY 4.0), distributed by arXiv under the Creative Commons Attribution 4.0 International licence, http://creativecommons.org/licenses/by/4.0/, as recorded in the arXiv licence metadata for this exact version and shown on the abstract page https://arxiv.org/abs/2309.15235v3. Full licence notice: \texttt{licenses/arxiv-2309.15235-CC-BY-4.0.txt}.\par\smallskip

\noindent\textit{Editorial scope.} Counted unit: the article's theorem t17 (source0.tex lines 2328--2380). The article's complete original proof of it is delivered with it (source0.tex lines 2382--2672, one \texttt{proof} environment of 291 lines). No mathematics is written, repaired or re-derived: the statement and the proof are the article's own lines, in the article's own notation, and no retained line is substituted or re-flowed.  Retained uncounted as the source-local context the proof needs: lines 1870--1891; lines 2072--2112; lines 2194--2216; lines 2220--2248.  Nothing was omitted from any retained span, so the package carries no omission record.  No retained line contains a citation, so the wrapper carries no bibliography and no bibliography entry.  No retained line contains a figure environment, an image-inclusion command or drawing code, so no image is delivered and the package declares no asset dependency.  Editorial formatting only, in addition: an AMS \texttt{amsart} preamble that restates the article's own declarations and macros used by the retained lines, namely the environments \texttt{proposition}, \texttt{definition}, \texttt{lemma}, \texttt{theorem} on separate counters, together with the standard packages those lines need.  The retained environments print in this document as Proposition 1, Definition 1, Lemma 1, Theorem 1 in order of appearance, whereas the article prints the same items with the numbering of its own full document.  The complete native source of the article as distributed in the arXiv e-print for this exact version is bundled unchanged as \texttt{sources/147600/source0.tex}.
\begin{proposition}[Stieltjes transform of the limit measure]
\label{prop:stieltjes-root}
For $w$ of sufficiently large modulus, let $\mathbf z_1(w,y)$ be the unique
solution of
 \[
U_y(z)=w
\]
that    satisfies $\mathbf z_1(w,y)\to1$ as $w\to\infty$. Then
\begin{equation}
\mathrm{St}_{\mathbf m_y}(w)
=
\operatorname{Log}\!\left(
\frac{\mathbf z_1(w,y)-1+s}{s}
\right),
\label{eq:st-root-large-w}
\end{equation}
where the logarithm is normalized by requiring the right-hand side to tend to
$0$ as $w\to\infty$.  Moreover, the identity extends by analytic
continuation to every connected domain that meets the large-$w$ domain and on
which the physical branch $\mathbf z_1(\,\cdot\,,y)$ and the normalized
logarithm are holomorphic.
\end{proposition}

\begin{proposition}[Local existence of the density]
\label{prop:density-existence}
Fix \(y\), and let \(I\subset\mathbb R\) be an open interval.
Let \(\mathbf z_1(w,y)\) be the solution branch introduced in
Proposition~\ref{prop:stieltjes-root}. For \(w\) in the large-\(w\)
domain, set
\[
    G_y(w)
    :=
    \operatorname{Log}\!\left(
        \frac{\mathbf z_1(w,y)-1+s}{s}
    \right),
\]
where the branch of \(\operatorname{Log}\) is the one used in
Proposition~\ref{prop:stieltjes-root}.

Assume that \(G_y\) admits a holomorphic continuation to an open
complex neighborhood \(\Omega\) of \(I\). Then, by
Proposition~\ref{prop:stieltjes-root} and the uniqueness of analytic
continuation,
\[
    \mathrm{St}_{\mathbf m_y}(w)=G_y(w),
    \qquad w\in\Omega\cap\mathbb H.
\]
Then the restriction of \(\mathbf m_y\) to \(I\) is absolutely
continuous and has the continuous density
\[
    f_y(\chi)
    =
    -\frac{1}{\pi}\Im G_y(\chi),
    \qquad \chi\in I.
\]
More precisely,
\[
    -\frac{1}{\pi}
    \Im\mathrm{St}_{\mathbf m_y}(\chi+\mathbf i\delta)
    \longrightarrow
    f_y(\chi)
\]
locally uniformly for \(\chi\in I\) as \(\delta\downarrow0\).
\end{proposition}

\begin{definition}\label{df28}
A point $(\chi,y)\in\mathcal R$ is called a \emph{regular liquid
point} if there is an open interval $I\ni\chi$ and a connected complex
neighborhood $\Omega$ of $I$ satisfying the assumptions of
Proposition~\ref{prop:density-existence}, and if the continuous density
supplied by that proposition satisfies
\[
    0<f_y(\chi)<1.
\]
Let
\[
    \mathcal L
    :=
    \{(\chi,y)\in\mathcal R:(\chi,y)\text{ is a regular liquid point}\}.
\]
By Proposition~\ref{prop:density-existence}, the limiting measure
$\mathbf m_y$ has a continuous density on every horizontal interval
contained in $\mathcal L$. The particle-spacing estimate gives only the
non-strict bound $0\leq f_y\leq1$; the strict inequalities above are part
of the definition. We call $\mathcal L$ the regular liquid region and
call its boundary $\partial\mathcal L$ the frozen boundary.

\end{definition}

\begin{lemma}[Nonreal roots in the regular liquid region]
\label{lem:liquid-nonreal-root}
Let \((\chi,y)\in\mathcal L\), and let
\(\mathbf z_1(\chi,y)\) denote the boundary value at \(w=\chi\) of the
solution branch introduced in
Proposition~\ref{prop:stieltjes-root}. Then
\[
    \Im \mathbf z_1(\chi,y)<0.
\]
Consequently,
\[
    \mathbf z_+(\chi,y)
    :=
    \overline{\mathbf z_1(\chi,y)}
\]
lies in the upper half-plane and satisfies
\[
    U_y(\mathbf z_+(\chi,y))=\chi.
\]
If the equation \(U_y(z)=\chi\) has at most one nonreal
complex-conjugate pair, then
\[
    \bigl\{
        \mathbf z_1(\chi,y),
        \mathbf z_+(\chi,y)
    \bigr\}
\]
is that unique pair.
\end{lemma}

\begin{theorem}[Multiple-root condition at a horizontal frozen boundary]
\label{t17}
Fix \(y\), and set
\[
    \mathcal L_y
    :=
    \{x\in\mathbb R:(x,y)\in\mathcal L\},
    \qquad
    \mathcal R_y
    :=
    \{x\in\mathbb R:(x,y)\in\mathcal R\}.
\]
Let \(I\) be a connected component of \(\mathcal L_y\), and let
\(\chi\) be a finite endpoint of \(I\) lying in
\(\operatorname{int}(\mathcal R_y)\). Choose a sequence
\[
    \chi_k\in I,
    \qquad
    \chi_k\longrightarrow\chi,
\]
and set
\[
    z_k:=\mathbf z_+(\chi_k,y)\in\mathbb H,
\]
where \(\mathbf z_+\) is defined in
Lemma~\ref{lem:liquid-nonreal-root}.

Assume that, after passing to a subsequence,
\[
    z_k\longrightarrow z_*
\]
for some finite \(z_*\in\mathbb C\), and that \(U_y\) is holomorphic on
a conjugation-invariant open neighborhood of
\(\{z_*,\overline{z_*}\}\). Assume also that, for every real \(x\)
sufficiently close to \(\chi\), the equation
\[
    U_y(z)=x
\]
has at most one nonreal complex-conjugate pair of roots.

Then \(z_*\in\mathbb R\) and
\[
    U_y(z_*)=\chi,
    \qquad
    U_y'(z_*)=0.
\]
Thus \(U_y(z)-\chi\) has a real zero of multiplicity at least two. If,
in addition,
\[
    U_y''(z_*)\neq0,
\]
then the multiplicity is exactly two.
\end{theorem}

\begin{proof}
For every \(k\),
\[
    U_y(z_k)=\chi_k.
\]
Since \(z_k\to z_*\), \(\chi_k\to\chi\), and \(U_y\) is holomorphic near
\(z_*\), we obtain
\[
    U_y(z_*)=\chi.
\]
Moreover, \(z_k\in\mathbb H\), so \(z_*\) lies in the closed upper
half-plane. We first prove that \(z_*\) is real.

Suppose, to the contrary, that \(z_*\in\mathbb H\). We distinguish two
cases.

\medskip
\noindent\emph{Case 1: \(U_y'(z_*)\neq0\).}
By the complex inverse-function theorem, there are a disk \(V\) centered at
\(\chi\), a neighborhood \(W_+\) of \(z_*\), and a holomorphic inverse branch
\[
    \zeta_+:V\longrightarrow W_+
\]
such that
\[
    U_y(\zeta_+(w))=w,
    \qquad
    \zeta_+(\chi)=z_*.
\]
After shrinking \(V\), we may assume that \(V\) is invariant under complex
conjugation. By the Schwarz symmetry assumed in the statement, the function
\[
    \zeta_-(w):=\overline{\zeta_+(\overline w)},
    \qquad w\in V,
\]
is holomorphic and satisfies
\[
    U_y(\zeta_-(w))=w,
    \qquad
    \zeta_-(\chi)=\overline{z_*}.
\]
Shrinking \(V\) once more, we may assume that
\[
    \zeta_+(V\cap\mathbb R)\subset\mathbb H,
    \qquad
    \zeta_-(V\cap\mathbb R)
    \subset\{z\in\mathbb C:\Im z<0\},
\]
and that
\[
    \frac{\zeta_-(w)-1+s}{s}\neq0,
    \qquad w\in V.
\]
For all sufficiently large \(k\), local uniqueness of the inverse branches
gives
\[
    \zeta_+(\chi_k)=z_k,
    \qquad
    \zeta_-(\chi_k)=\overline{z_k}
    =\mathbf z_1(\chi_k,y),
\]
where the last equality follows from
Lemma~\ref{lem:liquid-nonreal-root}.

Fix one such index \(k_0\). Since \((\chi_{k_0},y)\) is a regular liquid
point, Definition~\ref{df28} provides an open interval
\[
    J_{k_0}\ni\chi_{k_0}
\]
and a connected open complex neighborhood \(\Omega_{k_0}\) of \(J_{k_0}\)
to which the logarithmic branch
\[
    G_y(w)
    =
    \operatorname{Log}\!\left(
        \frac{\mathbf z_1(w,y)-1+s}{s}
    \right),
\]
initially defined for sufficiently large \(\lvert w\rvert\), admits a
single-valued holomorphic continuation. Denote this continuation by
\[
    G_{y,k_0}:\Omega_{k_0}\longrightarrow\mathbb C.
\]
By Proposition~\ref{prop:stieltjes-root} and uniqueness of analytic
continuation,
\[
    G_{y,k_0}(w)=\mathrm{St}_{\mathbf m_y}(w),
    \qquad w\in\Omega_{k_0}\cap\mathbb H.
\]

Set
\[
    \widehat\zeta_{y,k_0}(w)
    :=
    1-s+s\exp\bigl(G_{y,k_0}(w)\bigr).
\]
After shrinking \(\Omega_{k_0}\) around \(\chi_{k_0}\), if necessary, the
image of \(\widehat\zeta_{y,k_0}\) lies in a domain on which \(U_y\) is
holomorphic. On a nonempty open subset of
\(\Omega_{k_0}\cap\mathbb H\), Proposition~\ref{prop:stieltjes-root}
shows that
\[
    U_y\bigl(\widehat\zeta_{y,k_0}(w)\bigr)=w.
\]
The identity theorem therefore gives the same equality throughout the
shrunken neighborhood. At \(w=\chi_{k_0}\),
\[
    \widehat\zeta_{y,k_0}(\chi_{k_0})
    =\mathbf z_1(\chi_{k_0},y)
    =\zeta_-(\chi_{k_0}).
\]
Both \(\widehat\zeta_{y,k_0}\) and \(\zeta_-\) are thus local inverse
branches of \(U_y\) taking the same value at \(\chi_{k_0}\). Since this root
is simple, local uniqueness implies that they agree on a connected open
neighborhood of \(\chi_{k_0}\) contained in
\(V\cap\Omega_{k_0}\).

Because \(V\) is simply connected and
\((\zeta_-(w)-1+s)/s\) does not vanish on \(V\), there exists a holomorphic
function
\[
    \widetilde G_y:V\longrightarrow\mathbb C
\]
such that
\[
    \exp\bigl(\widetilde G_y(w)\bigr)
    =
    \frac{\zeta_-(w)-1+s}{s},
    \qquad w\in V.
\]
Choose this logarithmic branch so that it agrees with \(G_{y,k_0}\) on the
open neighborhood of \(\chi_{k_0}\) constructed above. Hence
\(\widetilde G_y\) and \(\mathrm{St}_{\mathbf m_y}\) agree on a nonempty
open subset of \(V\cap\mathbb H\). Since both functions are holomorphic on
the connected set \(V\cap\mathbb H\), the identity theorem yields
\[
    \mathrm{St}_{\mathbf m_y}(w)=\widetilde G_y(w),
    \qquad w\in V\cap\mathbb H.
\]
Thus the Stieltjes-transform branch extends through \(\chi\) to a full
complex neighborhood, not merely from the liquid side.

Since \(\chi\in\operatorname{int}_{\mathbb R}(\mathcal R_y)\), choose an
open interval \(J_0\ni\chi\) such that
\[
    \overline{J_0}\subset V\cap\mathbb R,
    \qquad
    J_0\subset\mathcal R_y.
\]
Proposition~\ref{prop:density-existence}, applied to
\(\widetilde G_y\), shows that \(\mathbf m_y\) is absolutely continuous on
\(J_0\) with continuous density
\[
    f_y(x)
    =
    -\frac{1}{\pi}\Im\widetilde G_y(x),
    \qquad x\in J_0.
\]
For \(x\in I\cap J_0\), we have \((x,y)\in\mathcal L\), and hence
\[
    0<f_y(x)<1.
\]
Therefore,
\[
    -\pi<\Im\widetilde G_y(x)<0,
    \qquad x\in I\cap J_0.
\]
Since \(\chi\) is an endpoint of \(I\), points of \(I\cap J_0\) approach
\(\chi\) from the liquid side. By continuity,
\[
    -\pi\leq\Im\widetilde G_y(\chi)\leq0.
\]
On the other hand,
\[
    \exp\bigl(\widetilde G_y(\chi)\bigr)
    =
    \frac{\zeta_-(\chi)-1+s}{s}
    =
    \frac{\overline{z_*}-1+s}{s}
\]
is nonreal. Hence \(\Im\widetilde G_y(\chi)\) is neither \(0\) nor
\(-\pi\), and thus
\[
    -\pi<\Im\widetilde G_y(\chi)<0.
\]
By continuity, after shrinking to a smaller open interval
\(J\ni\chi\) with \(J\subset J_0\), we have
\[
    0<-\frac{1}{\pi}\Im\widetilde G_y(x)<1,
    \qquad x\in J.
\]
Moreover, \(\zeta_-\) is a holomorphic inverse branch on \(V\), so
\[
    U_y'(\zeta_-(w))\,\zeta_-'(w)=1
\]
and therefore \(U_y'(\zeta_-(w))\neq0\) on \(V\). Thus the same root and
logarithmic branches verify all the regularity conditions in
Definition~\ref{df28} for every \(x\in J\). Consequently,
\[
    J\subset\mathcal L_y.
\]
The interval \(J\) intersects \(I\) and contains the endpoint
\(\chi\notin I\). Hence \(I\cup J\) is a connected subset of
\(\mathcal L_y\) strictly larger than the connected component \(I\), a
contradiction.

\medskip
\noindent\emph{Case 2: \(U_y'(z_*)=0\).}
Let \(m\geq2\) be the multiplicity of \(z_*\) as a zero of
\[
    F(z):=U_y(z)-\chi.
\]
Choose a closed disk \(D_+\Subset\mathbb H\), centered at \(z_*\) and
contained in the holomorphy domain of \(U_y\), such that \(F\) has no zero
on \(\partial D_+\) and no zero in \(D_+\) other than \(z_*\), counted with
multiplicity \(m\). Set
\[
    c:=\min_{z\in\partial D_+}|F(z)|>0.
\]
For every real \(\varepsilon\) with \(|\varepsilon|<c\), Rouch\'e's theorem
shows that
\[
    U_y(z)-(\chi+\varepsilon)=F(z)-\varepsilon
\]
has exactly \(m\) zeros in \(D_+\), counted with multiplicity.

Since \(U_y\) is nonconstant, \(U_y'\) is not identically zero and has only
finitely many zeros in \(\overline{D_+}\). Hence the set
\[
    \mathcal C
    :=
    \bigl\{
        U_y(c_0)-\chi:
        c_0\in\overline{D_+},\ U_y'(c_0)=0
    \bigr\}
    \cap\mathbb R
\]
is finite. We may therefore choose a nonzero real \(\varepsilon\),
arbitrarily small and within the neighborhood in which the one-pair
assumption holds, such that
\[
    |\varepsilon|<c,
    \qquad
    \varepsilon\notin\mathcal C.
\]
For this choice, all \(m\) zeros in \(D_+\) are simple and hence distinct.
Since \(m\geq2\), the equation
\[
    U_y(z)=\chi+\varepsilon
\]
has at least two distinct roots in the upper half-plane. By Schwarz
symmetry, their conjugates produce at least two nonreal complex-conjugate
pairs, contradicting the one-pair assumption.

The two cases exclude \(z_*\in\mathbb H\). Since \(z_*\) lies in the closed
upper half-plane, we conclude that
\[
    z_*\in\mathbb R.
\]

It remains to prove that the limiting real root is multiple. Since
\(z_*\in\mathbb R\),
\[
    \overline{z_k}\longrightarrow z_*.
\]
For each \(k\), the two distinct points \(z_k\) and \(\overline{z_k}\) are
zeros of
\[
    U_y(z)-\chi_k.
\]
Suppose that \(z_*\) were a simple zero of \(U_y(z)-\chi\). Choose a closed
disk \(D\), centered at \(z_*\) and contained in the holomorphy domain of
\(U_y\), such that \(U_y(z)-\chi\) has no zero on \(\partial D\) and exactly
one zero in \(D\), counted with multiplicity. Set
\[
    c_D:=\min_{z\in\partial D}|U_y(z)-\chi|>0.
\]
For all sufficiently large \(k\),
\[
    |\chi_k-\chi|<c_D.
\]
Rouch\'e's theorem then implies that \(U_y(z)-\chi_k\) has exactly one zero
in \(D\), counted with multiplicity. However, for all sufficiently large
\(k\), both \(z_k\) and \(\overline{z_k}\) lie in \(D\), and they are
distinct. This contradiction shows that \(z_*\) is a zero of
\(U_y(z)-\chi\) of multiplicity at least two. Therefore,
\[
    U_y'(z_*)=0.
\]
If \(U_y''(z_*)\neq0\), then the multiplicity is exactly two.
\end{proof}

\end{document}
